% =========================================================
% Chapter 10: Fluid Mechanics
% POSN 1 Physics Preparation Notes
% Language: Thai
% Compiler: XeLaTeX
% =========================================================

\chapter{กลศาสตร์ของไหล}

\begin{center}
    {\Large \textbf{Fluid Mechanics}}\\[4pt]
    {\normalsize เอกสารเตรียมสอบคัดเลือก สอวน. ฟิสิกส์ ค่าย 1}\\[6pt]
    {\normalsize จัดทำโดย \textbf{Tames Thanakrit Damduan}}\\[2pt]
    {\footnotesize \textcopyright\ 2026 Tames Thanakrit Damduan. All rights reserved.}\\[8pt]
\end{center}

\section*{เป้าหมายของบทเรียน}

หลังจากเรียนบทนี้ นักเรียนควรทำสิ่งต่อไปนี้ได้

\begin{enumerate}
    \item เข้าใจความหมายของความหนาแน่น ความดัน และความดันเกจ
    \item ใช้ความสัมพันธ์ความดันในของเหลวอยู่นิ่งได้
    \item วิเคราะห์หลอดปรอท แมนอมิเตอร์ และหลอดรูปตัวยูได้
    \item ใช้หลักของอาร์คิมีดีสเพื่อหาแรงพยุงได้
    \item วิเคราะห์วัตถุลอย จม และลอยนิ่งในของไหลได้
    \item ใช้สมการความต่อเนื่องของการไหลได้
    \item ใช้สมการแบร์นูลลีและกฎของตอร์ริเชลลีได้
    \item เข้าใจแนวคิดของความหนืด แรงต้าน และความตึงผิวในระดับที่ใช้กับข้อสอบได้
\end{enumerate}

\begin{tcolorbox}[title=แนวคิดสำคัญของบทนี้]
ของไหลต่างจากวัตถุแข็งเกร็ง เพราะของไหลเปลี่ยนรูปร่างได้ง่ายและส่งแรงผ่านความดัน จุดสำคัญของบทนี้คือการแยกให้ชัดว่าโจทย์เป็นของไหลอยู่นิ่งหรือของไหลกำลังไหล ถ้าอยู่นิ่งให้คิดจากความดันและแรงพยุง ถ้ากำลังไหลให้คิดจากความต่อเนื่องและพลังงานของของไหล
\end{tcolorbox}

\newpage{}

% =========================================================
\section{ความหนาแน่นและความดัน}

\subsection{ความหนาแน่น}

ความหนาแน่นบอกว่ามวลถูกบรรจุอยู่ในปริมาตรมากน้อยเพียงใด

ถ้าวัตถุมีมวล

\[
m
\]

และมีปริมาตร

\[
V
\]

นิยามความหนาแน่นเป็น

\[
\rho=\frac{m}{V}
\]

\begin{tcolorbox}[title=นิยามความหนาแน่น]
\[
\rho=\frac{m}{V}
\]

หน่วย SI คือ

\[
\si{kg/m^3}
\]
\end{tcolorbox}

ถ้าวัตถุมีความหนาแน่นมาก แปลว่ามีมวลมากต่อปริมาตรหนึ่งหน่วย

ตัวอย่างเช่น น้ำมีความหนาแน่นประมาณ

\[
\rho_{\text{water}}=1000\,\si{kg/m^3}
\]

ส่วนอากาศมีความหนาแน่นน้อยกว่าน้ำมาก

\subsection{ความดัน}

ความดันเกิดจากแรงที่กระทำตั้งฉากกับพื้นที่

ถ้าแรงตั้งฉากมีขนาด

\[
F
\]

กระทำบนพื้นที่

\[
A
\]

นิยามความดันเป็น

\[
P=\frac{F}{A}
\]

\begin{tcolorbox}[title=นิยามความดัน]
\[
P=\frac{F}{A}
\]

หน่วย SI คือ

\[
\si{N/m^2}
\]

หรือเรียกว่า

\[
\si{Pa}
\]
\end{tcolorbox}

ถ้าแรงเท่าเดิมแต่พื้นที่เล็กลง ความดันจะมากขึ้น เช่น ปลายเข็มมีพื้นที่เล็กมากจึงเกิดความดันสูง

\subsection{ความดันสัมบูรณ์และความดันเกจ}

ความดันสัมบูรณ์คือความดันจริงที่วัดเทียบกับสุญญากาศ

\[
P_{\text{abs}}
\]

ความดันเกจคือความดันที่วัดเทียบกับความดันบรรยากาศ

\[
P_{\text{gauge}}
\]

ดังนั้น

\[
P_{\text{abs}}=P_{\text{atm}}+P_{\text{gauge}}
\]

\begin{tcolorbox}[title=ความดันสัมบูรณ์และความดันเกจ]
\[
P_{\text{abs}}=P_{\text{atm}}+P_{\text{gauge}}
\]

ถ้าโจทย์บอกว่า ``ความดันเกจ'' ต้องบวกความดันบรรยากาศก่อนเมื่อต้องการความดันสัมบูรณ์
\end{tcolorbox}

\begin{examplebox}{ตัวอย่าง: ความดันจากแรง}
แรงขนาด $200\,\si{N}$ กระทำตั้งฉากบนพื้นที่ $0.50\,\si{m^2}$ จงหาความดัน

เริ่มจากนิยาม

\[
P=\frac{F}{A}
\]

แทนค่า

\[
P=\frac{200}{0.50}
\]

\[
P=400\,\si{Pa}
\]
\end{examplebox}

\begin{exercisebox}{แบบฝึกหัด 10.1}
\begin{enumerate}
    \item วัตถุมวล $2.0\,\si{kg}$ มีปริมาตร $4.0\times 10^{-3}\,\si{m^3}$ จงหาความหนาแน่น
    \item แรง $500\,\si{N}$ กระทำตั้งฉากบนพื้นที่ $0.25\,\si{m^2}$ จงหาความดัน
    \item ถ้าความดันเกจเท่ากับ $2.0\times 10^5\,\si{Pa}$ และความดันบรรยากาศเท่ากับ $1.0\times 10^5\,\si{Pa}$ จงหาความดันสัมบูรณ์
    \item เหตุใดรองเท้าส้นเข็มจึงกดพื้นด้วยความดันสูงกว่ารองเท้าพื้นกว้าง ถ้าแรงกดเท่ากัน
\end{enumerate}
\end{exercisebox}

% =========================================================
\section{ความดันในของเหลวอยู่นิ่ง}

\subsection{ความดันเพิ่มตามความลึก}

พิจารณาของเหลวที่อยู่นิ่ง ความดันที่ระดับลึกกว่าจะมากกว่าความดันที่ระดับตื้นกว่า เพราะของเหลวด้านล่างต้องรับน้ำหนักของของเหลวที่อยู่เหนือมัน

พิจารณาชั้นของเหลวที่มีพื้นที่หน้าตัด

\[
A
\]

และมีความสูง

\[
h
\]

ปริมาตรของชั้นของเหลวคือ

\[
V=Ah
\]

มวลของชั้นของเหลวคือ

\[
m=\rho V=\rho Ah
\]

น้ำหนักของชั้นของเหลวคือ

\[
mg=\rho Ahg
\]

แรงนี้กระจายบนพื้นที่ $A$ ดังนั้นความดันที่เพิ่มขึ้นจากน้ำหนักของของเหลวคือ

\[
\Delta P=\frac{\rho Ahg}{A}
\]

จึงได้

\[
\Delta P=\rho gh
\]

ถ้าผิวบนของของเหลวมีความดัน

\[
P_0
\]

ความดันที่ลึกลงไป $h$ คือ

\[
P=P_0+\rho gh
\]

\begin{tcolorbox}[title=ความดันในของเหลวอยู่นิ่ง]
\[
P=P_0+\rho gh
\]

ถ้าผิวของเหลวเปิดสู่อากาศ

\[
P_0=P_{\text{atm}}
\]
\end{tcolorbox}

\subsection{ความดันขึ้นกับความลึก ไม่ขึ้นกับรูปภาชนะ}

จากสูตร

\[
P=P_0+\rho gh
\]

ความดันขึ้นกับ

\[
\rho,\quad g,\quad h
\]

แต่ไม่ขึ้นกับรูปร่างของภาชนะ

ดังนั้นจุดสองจุดในของเหลวชนิดเดียวกันที่อยู่นิ่งและอยู่ลึกเท่ากัน จะมีความดันเท่ากัน

\begin{tcolorbox}[title=หลักสำคัญของของเหลวอยู่นิ่ง]
ในของเหลวชนิดเดียวกันที่อยู่นิ่ง จุดที่อยู่ระดับเดียวกันมีความดันเท่ากัน
\end{tcolorbox}

\subsection{หลักของปาสคาล}

เมื่อเพิ่มความดันให้ของไหลที่ถูกปิดล้อม ความดันที่เพิ่มขึ้นจะส่งผ่านไปทุกจุดในของไหล

ถ้าลูกสูบเล็กมีพื้นที่

\[
A_1
\]

และถูกกดด้วยแรง

\[
F_1
\]

ความดันที่สร้างคือ

\[
P=\frac{F_1}{A_1}
\]

ความดันนี้ส่งไปยังลูกสูบใหญ่พื้นที่

\[
A_2
\]

จึงเกิดแรง

\[
F_2=PA_2
\]

ดังนั้น

\[
\frac{F_1}{A_1}=\frac{F_2}{A_2}
\]

\begin{tcolorbox}[title=หลักของปาสคาล]
\[
\frac{F_1}{A_1}=\frac{F_2}{A_2}
\]

ใช้กับระบบไฮดรอลิก เช่น แม่แรงและเบรกไฮดรอลิก
\end{tcolorbox}

\begin{examplebox}{ตัวอย่าง: ความดันใต้น้ำ}
จุดหนึ่งอยู่ลึกจากผิวน้ำ $5.0\,\si{m}$ จงหาความดันเกจที่จุดนั้น ใช้ $\rho=1000\,\si{kg/m^3}$ และ $g=10\,\si{m/s^2}$

ความดันเกจคือส่วนที่เพิ่มจากความดันบรรยากาศ

\[
P_{\text{gauge}}=\rho gh
\]

แทนค่า

\[
P_{\text{gauge}}=(1000)(10)(5.0)
\]

\[
P_{\text{gauge}}=5.0\times 10^4\,\si{Pa}
\]
\end{examplebox}

\begin{exercisebox}{แบบฝึกหัด 10.2}
\begin{enumerate}
    \item จงหาความดันเกจที่ความลึก $2.0\,\si{m}$ ในน้ำ ใช้ $g=10\,\si{m/s^2}$
    \item ถ้าความดันที่ก้นถังเปิดฝาเท่ากับ $1.4\times 10^5\,\si{Pa}$ และความดันบรรยากาศเท่ากับ $1.0\times 10^5\,\si{Pa}$ จงหาความลึกของน้ำ
    \item ลูกสูบเล็กมีพื้นที่ $2\,\si{cm^2}$ และลูกสูบใหญ่มีพื้นที่ $50\,\si{cm^2}$ ถ้าออกแรงที่ลูกสูบเล็ก $20\,\si{N}$ จงหาแรงที่ลูกสูบใหญ่
    \item จุดสองจุดอยู่ในน้ำที่ระดับลึกเท่ากัน แต่ภาชนะมีรูปร่างต่างกัน ความดันที่จุดทั้งสองเท่ากันหรือไม่ เพราะเหตุใด
\end{enumerate}
\end{exercisebox}

% =========================================================
\section{การวัดความดัน}

\subsection{บารอมิเตอร์ปรอท}

บารอมิเตอร์ใช้วัดความดันบรรยากาศ โดยเทียบความดันบรรยากาศกับความดันจากคอลัมน์ปรอท

ถ้าความสูงของปรอทคือ

\[
h
\]

จะได้

\[
P_{\text{atm}}=\rho_{\text{Hg}}gh
\]

ในหน่วยมิลลิเมตรปรอท เขียนได้ว่า

\[
1\,\si{atm}\approx 760\,\si{mmHg}
\]

\begin{tcolorbox}[title=บารอมิเตอร์]
\[
P_{\text{atm}}=\rho_{\text{Hg}}gh
\]

หรือประมาณ

\[
P_{\text{atm}}=760\,\si{mmHg}
\]
\end{tcolorbox}

\subsection{แมนอมิเตอร์และหลอดรูปตัวยู}

หลักของแมนอมิเตอร์คือ จุดในของเหลวชนิดเดียวกันที่อยู่ระดับเดียวกันมีความดันเท่ากัน

ถ้าของเหลวในหลอดรูปตัวยูมีความหนาแน่น $\rho$ และระดับต่างกัน $h$ ความดันต่างกันคือ

\[
\Delta P=\rho gh
\]

\begin{tcolorbox}[title=หลักการแก้โจทย์หลอดรูปตัวยู]
\begin{enumerate}
    \item เลือกระดับอ้างอิงแนวนอนในของเหลวชนิดเดียวกัน
    \item เขียนความดันด้านซ้ายและด้านขวาที่ระดับเดียวกัน
    \item ใช้หลักว่าความดันเท่ากันที่ระดับเดียวกัน
    \item ระวังความหนาแน่นของของเหลวแต่ละชั้น
\end{enumerate}
\end{tcolorbox}

\begin{examplebox}{ตัวอย่าง: ความดันในหน่วยมิลลิเมตรปรอท}
ความดันบรรยากาศเท่ากับ $750\,\si{mmHg}$ ถ้าในหลอดแก้วมีผิวปรอทสูงกว่าผิวปรอทในอ่าง $10\,\si{mm}$ จงหาความดันเหนือผิวปรอทในหลอด

ความดันที่ระดับเดียวกันในปรอทต้องเท่ากัน

ด้านนอกมีความดันบรรยากาศ

\[
P_{\text{atm}}=750\,\si{mmHg}
\]

ด้านในมีความดันเหนือผิวปรอทในหลอดเป็น $P$ และมีคอลัมน์ปรอทสูงกว่า $10\,\si{mmHg}$

ดังนั้น

\[
P+10=750
\]

จึงได้

\[
P=740\,\si{mmHg}
\]
\end{examplebox}

% =========================================================
\section{แรงพยุงและหลักของอาร์คิมีดีส}

\subsection{แรงพยุงเกิดจากอะไร}

เมื่อวัตถุจมอยู่ในของเหลว ความดันที่ด้านล่างมากกว่าความดันที่ด้านบน เพราะด้านล่างอยู่ลึกกว่า

ดังนั้นแรงจากของเหลวที่ดันขึ้นด้านล่างจะมากกว่าแรงที่กดลงด้านบน ผลรวมจึงเป็นแรงลัพธ์ขึ้นด้านบน

แรงนี้เรียกว่า

\[
\text{แรงพยุง}
\]

\subsection{หลักของอาร์คิมีดีส}

หลักของอาร์คิมีดีสกล่าวว่า แรงพยุงที่ของไหลกระทำต่อวัตถุมีขนาดเท่ากับน้ำหนักของของไหลที่ถูกแทนที่

ถ้าวัตถุแทนที่ของไหลปริมาตร

\[
V_{\text{disp}}
\]

ของไหลมีความหนาแน่น

\[
\rho_f
\]

มวลของของไหลที่ถูกแทนที่คือ

\[
\rho_f V_{\text{disp}}
\]

ดังนั้นน้ำหนักของของไหลที่ถูกแทนที่คือ

\[
\rho_f g V_{\text{disp}}
\]

แรงพยุงจึงมีขนาด

\[
B=\rho_f gV_{\text{disp}}
\]

\begin{tcolorbox}[title=หลักของอาร์คิมีดีส]
\[
B=\rho_f gV_{\text{disp}}
\]

แรงพยุงมีทิศขึ้น และมีขนาดเท่ากับน้ำหนักของของไหลที่ถูกแทนที่
\end{tcolorbox}

\subsection{วัตถุจมทั้งหมด}

ถ้าวัตถุจมอยู่ในของไหลทั้งหมด

\[
V_{\text{disp}}=V_{\text{object}}
\]

ดังนั้น

\[
B=\rho_f gV_{\text{object}}
\]

น้ำหนักของวัตถุคือ

\[
W=\rho_{\text{object}}gV_{\text{object}}
\]

ถ้า

\[
\rho_{\text{object}}>\rho_f
\]

วัตถุจม

ถ้า

\[
\rho_{\text{object}}<\rho_f
\]

วัตถุลอยขึ้น

ถ้า

\[
\rho_{\text{object}}=\rho_f
\]

วัตถุลอยนิ่งในของไหล

\subsection{วัตถุลอยบางส่วน}

ถ้าวัตถุลอยนิ่งบนผิวของไหล แรงพยุงเท่ากับน้ำหนักของวัตถุ

\[
B=W
\]

ดังนั้น

\[
\rho_f gV_{\text{disp}}=\rho_{\text{object}}gV_{\text{object}}
\]

ตัด $g$ ออก

\[
\frac{V_{\text{disp}}}{V_{\text{object}}}
=
\frac{\rho_{\text{object}}}{\rho_f}
\]

\begin{tcolorbox}[title=เศษส่วนของวัตถุที่จมน้ำเมื่อวัตถุลอยนิ่ง]
\[
\frac{V_{\text{disp}}}{V_{\text{object}}}
=
\frac{\rho_{\text{object}}}{\rho_f}
\]

ใช้เฉพาะกรณีวัตถุลอยนิ่ง
\end{tcolorbox}

\begin{examplebox}{ตัวอย่าง: ก้อนน้ำแข็งลอยในน้ำ}
ก้อนน้ำแข็งมีความหนาแน่น

\[
\rho_{\text{ice}}=0.92\,\si{g/cm^3}
\]

ลอยอยู่ในน้ำที่มีความหนาแน่น

\[
\rho_{\text{water}}=1.00\,\si{g/cm^3}
\]

จงหาสัดส่วนของปริมาตรน้ำแข็งที่จมอยู่ในน้ำ

ใช้

\[
\frac{V_{\text{disp}}}{V_{\text{object}}}
=
\frac{\rho_{\text{ice}}}{\rho_{\text{water}}}
\]

แทนค่า

\[
\frac{V_{\text{disp}}}{V_{\text{object}}}
=
\frac{0.92}{1.00}
\]

ดังนั้น

\[
\frac{V_{\text{disp}}}{V_{\text{object}}}=0.92
\]

แปลว่า น้ำแข็งจมอยู่ใต้น้ำประมาณ $92\%$ ของปริมาตรทั้งหมด
\end{examplebox}

\begin{exercisebox}{แบบฝึกหัด 10.3}
\begin{enumerate}
    \item วัตถุปริมาตร $2.0\times 10^{-3}\,\si{m^3}$ จมอยู่ในน้ำทั้งหมด จงหาแรงพยุง
    \item วัตถุมวล $3.0\,\si{kg}$ และปริมาตร $2.0\times 10^{-3}\,\si{m^3}$ เมื่อนำไปจมน้ำทั้งหมด แรงลัพธ์ชี้ขึ้นหรือลง
    \item ไม้ความหนาแน่น $600\,\si{kg/m^3}$ ลอยในน้ำ จงหาสัดส่วนปริมาตรที่จมน้ำ
    \item วัตถุหนึ่งมีปริมาตร $V$ และความหนาแน่นครึ่งหนึ่งของน้ำ จงบอกว่าจมน้ำกี่ส่วนของปริมาตรทั้งหมด
\end{enumerate}
\end{exercisebox}

% =========================================================
\section{ของไหลอุดมคติและสมการความต่อเนื่อง}

\subsection{แบบจำลองของไหลอุดมคติ}

ในการวิเคราะห์ของไหลที่กำลังไหล เรามักเริ่มจากแบบจำลองของไหลอุดมคติ ซึ่งมีสมบัติดังนี้

\begin{enumerate}
    \item ไม่หนืด
    \item ไหลแบบสม่ำเสมอหรือไหลเป็นชั้น
    \item อัดตัวไม่ได้
    \item ไม่มีการหมุนวนในตัวของไหล
\end{enumerate}

นี่เป็นแบบจำลองที่ทำให้สมการง่ายขึ้น และเพียงพอสำหรับโจทย์พื้นฐานส่วนใหญ่

\subsection{อัตราการไหล}

ถ้าของไหลผ่านพื้นที่หน้าตัด

\[
A
\]

ด้วยอัตราเร็ว

\[
v
\]

ในเวลา $\Delta t$ ของไหลเคลื่อนที่ได้ระยะ

\[
v\Delta t
\]

ปริมาตรที่ผ่านหน้าตัดคือ

\[
\Delta V=Av\Delta t
\]

ดังนั้นอัตราการไหลเชิงปริมาตรคือ

\[
Q=\frac{\Delta V}{\Delta t}=Av
\]

\begin{tcolorbox}[title=อัตราการไหลเชิงปริมาตร]
\[
Q=Av
\]

หน่วยคือ

\[
\si{m^3/s}
\]
\end{tcolorbox}

\subsection{สมการความต่อเนื่อง}

ถ้าของไหลอัดตัวไม่ได้ ปริมาตรที่ไหลผ่านแต่ละหน้าตัดในเวลาเดียวกันต้องเท่ากัน

ดังนั้น

\[
A_1v_1=A_2v_2
\]

\begin{tcolorbox}[title=สมการความต่อเนื่อง]
\[
A_1v_1=A_2v_2
\]

ถ้าท่อแคบลง อัตราเร็วของของไหลจะเพิ่มขึ้น
\end{tcolorbox}

\begin{examplebox}{ตัวอย่าง: น้ำไหลในท่อ}
น้ำไหลในท่อที่พื้นที่หน้าตัด $4\,\si{cm^2}$ ด้วยอัตราเร็ว $2\,\si{m/s}$ ถ้าท่อลดพื้นที่เหลือ $1\,\si{cm^2}$ จงหาอัตราเร็วในส่วนแคบ

ใช้

\[
A_1v_1=A_2v_2
\]

แทนค่า

\[
(4)(2)=(1)v_2
\]

ดังนั้น

\[
v_2=8\,\si{m/s}
\]
\end{examplebox}

% =========================================================
\section{สมการแบร์นูลลี}

\subsection{แนวคิดพลังงานของของไหล}

สมการแบร์นูลลีเป็นการอนุรักษ์พลังงานต่อปริมาตรของของไหลอุดมคติที่ไหลตามเส้นกระแสเดียวกัน

พลังงานต่อปริมาตรมีสามส่วน

\[
\text{ความดัน }P
\]

\[
\text{พลังงานจลน์ต่อปริมาตร } \frac{1}{2}\rho v^2
\]

\[
\text{พลังงานศักย์โน้มถ่วงต่อปริมาตร } \rho gy
\]

ดังนั้นตามเส้นกระแสเดียวกัน

\[
P+\frac{1}{2}\rho v^2+\rho gy=\text{constant}
\]

\begin{tcolorbox}[title=สมการแบร์นูลลี]
\[
P+\frac{1}{2}\rho v^2+\rho gy=\text{constant}
\]

หรือระหว่างสองจุด

\[
P_1+\frac{1}{2}\rho v_1^2+\rho gy_1
=
P_2+\frac{1}{2}\rho v_2^2+\rho gy_2
\]
\end{tcolorbox}

\subsection{ข้อควรระวังในการใช้สมการแบร์นูลลี}

สมการแบร์นูลลีใช้ได้ดีเมื่อ

\begin{enumerate}
    \item ของไหลอัดตัวไม่ได้
    \item ความหนืดน้อย
    \item การไหลสม่ำเสมอ
    \item จุดที่พิจารณาอยู่บนเส้นกระแสเดียวกัน
\end{enumerate}

\begin{tcolorbox}[title=ข้อควรจำ]
แบร์นูลลีไม่ใช่สูตรความดันในของเหลวอยู่นิ่งโดยตรง แต่ถ้าของไหลอยู่นิ่ง คือ $v_1=v_2=0$ สมการแบร์นูลลีจะกลับไปเป็นสูตรความดันตามความลึก
\end{tcolorbox}

\subsection{กฎของตอร์ริเชลลี}

พิจารณาถังเปิดที่มีรูเล็กที่ความลึก $h$ ใต้ผิวน้ำ

จุดที่ผิวน้ำมีความดัน

\[
P_1=P_{\text{atm}}
\]

จุดที่รูเปิดก็มีความดัน

\[
P_2=P_{\text{atm}}
\]

ถ้าถังกว้างมาก ความเร็วของผิวน้ำเล็กมาก

\[
v_1\approx 0
\]

เลือกให้รูอยู่ที่ระดับ $y=0$ และผิวน้ำอยู่ที่ $y=h$

ใช้แบร์นูลลี

\[
P_{\text{atm}}+\rho gh
=
P_{\text{atm}}+\frac{1}{2}\rho v^2
\]

ตัด $P_{\text{atm}}$ และ $\rho$ ออก

\[
gh=\frac{1}{2}v^2
\]

จึงได้

\[
v=\sqrt{2gh}
\]

\begin{tcolorbox}[title=กฎของตอร์ริเชลลี]
\[
v=\sqrt{2gh}
\]

ใช้กับรูเล็กในถังเปิด เมื่อความเร็วของผิวน้ำเล็กมาก
\end{tcolorbox}

\subsection{กรณีพื้นที่ผิวน้ำไม่ใหญ่มาก}

ถ้าพื้นที่หน้าตัดของถังเป็น

\[
A
\]

และพื้นที่รูเป็น

\[
a
\]

จากสมการความต่อเนื่อง

\[
Av_{\text{surface}}=av_{\text{hole}}
\]

ดังนั้น

\[
v_{\text{surface}}=\frac{a}{A}v_{\text{hole}}
\]

ใช้แบร์นูลลีจากผิวน้ำถึงรู

\[
\rho gh+\frac{1}{2}\rho v_{\text{surface}}^2
=
\frac{1}{2}\rho v_{\text{hole}}^2
\]

แทน

\[
v_{\text{surface}}=\frac{a}{A}v_{\text{hole}}
\]

จะได้

\[
gh+\frac{1}{2}\left(\frac{a}{A}\right)^2v_{\text{hole}}^2
=
\frac{1}{2}v_{\text{hole}}^2
\]

จัดรูป

\[
2gh=
v_{\text{hole}}^2
\left[
1-\left(\frac{a}{A}\right)^2
\right]
\]

ดังนั้น

\[
v_{\text{hole}}
=
\sqrt{
\frac{2gh}
{1-\left(\frac{a}{A}\right)^2}
}
\]

\begin{tcolorbox}[title=รูที่ก้นถังเมื่อพื้นที่ถังไม่ใหญ่มาก]
\[
v_{\text{hole}}
=
\sqrt{
\frac{2gh}
{1-\left(\frac{a}{A}\right)^2}
}
\]

ถ้า $a\ll A$ สูตรนี้จะลดรูปเป็น

\[
v=\sqrt{2gh}
\]
\end{tcolorbox}

\begin{exercisebox}{แบบฝึกหัด 10.4}
\begin{enumerate}
    \item น้ำไหลผ่านท่อหน้าตัด $6\,\si{cm^2}$ ด้วยอัตราเร็ว $2\,\si{m/s}$ แล้วไหลเข้าสู่ท่อหน้าตัด $3\,\si{cm^2}$ จงหาอัตราเร็วในท่อส่วนหลัง
    \item ถังเปิดมีรูเล็กที่ลึกจากผิวน้ำ $1.25\,\si{m}$ ใช้ $g=10\,\si{m/s^2}$ จงหาอัตราเร็วของน้ำที่พุ่งออกจากรู
    \item น้ำไหลในท่อแนวนอน ถ้าท่อแคบลงและอัตราเร็วเพิ่มขึ้น ความดันเพิ่มหรือลด
    \item จงอธิบายว่าทำไมสมการแบร์นูลลีจึงเกี่ยวข้องกับการอนุรักษ์พลังงาน
\end{enumerate}
\end{exercisebox}

% =========================================================
\section{ความหนืด แรงต้าน และความตึงผิว}

\subsection{ความหนืด}

ความหนืดคือความต้านทานต่อการไหลของของไหล ของไหลที่หนืดมาก เช่น น้ำผึ้ง ไหลยากกว่าน้ำ

ในของไหลจริง พลังงานกลบางส่วนจะสูญเสียเป็นความร้อนเพราะความหนืด ดังนั้นสมการแบร์นูลลีแบบอุดมคติอาจใช้ไม่ได้ตรง ๆ ในทุกสถานการณ์

\subsection{แรงต้านในของไหล}

เมื่อวัตถุเคลื่อนที่ในของไหล มักมีแรงต้านทิศตรงข้ามกับการเคลื่อนที่

โจทย์บางข้อกำหนดแรงต้านในรูป

\[
F_D=Cv
\]

หรือ

\[
F_D=kv^2A
\]

ถ้าโจทย์ให้รูปของแรงต้านมา ให้ใช้รูปนั้นตามโจทย์ และตั้งสมการแรงตามแนวการเคลื่อนที่

\begin{tcolorbox}[title=แรงต้านกับความเร็วคงตัว]
ถ้าวัตถุเคลื่อนที่ในของไหลจนมีความเร็วคงตัว จะมี

\[
\sum F=0
\]

ตามแนวการเคลื่อนที่
\end{tcolorbox}

\subsection{ความตึงผิว}

ความตึงผิวเป็นแรงต่อความยาวบริเวณผิวของของเหลว ใช้สัญลักษณ์

\[
\gamma
\]

หน่วยคือ

\[
\si{N/m}
\]

สำหรับหยดของเหลวที่มีผิวเดียว ความดันด้านในมากกว่าด้านนอกเป็น

\[
\Delta P=\frac{2\gamma}{R}
\]

สำหรับฟองสบู่ที่มีผิวสองด้าน ความดันด้านในมากกว่าด้านนอกเป็น

\[
\Delta P=\frac{4\gamma}{R}
\]

\begin{tcolorbox}[title=ความดันเกินจากความตึงผิว]
หยดของเหลวผิวเดียว

\[
P_i-P_o=\frac{2\gamma}{R}
\]

ฟองสบู่สองผิว

\[
P_i-P_o=\frac{4\gamma}{R}
\]
\end{tcolorbox}

\subsection{การไหลหนืดในท่อ}

สำหรับของไหลหนืดที่ไหลแบบเป็นชั้นในท่อยาว รัศมีท่อมีผลมากต่ออัตราการไหล

ผลสำคัญคือ

\[
Q\propto R^4
\]

ดังนั้นถ้ารัศมีท่อลดลงเล็กน้อย อัตราการไหลอาจลดลงมาก

\begin{tcolorbox}[title=แนวคิดจากการไหลหนืด]
สำหรับการไหลหนืดในท่อ

\[
Q\propto R^4
\]

จุดนี้มักใช้เชิงแนวคิดมากกว่าการคำนวณเต็มในข้อสอบพื้นฐาน
\end{tcolorbox}

% =========================================================
\section{วิธีเลือกหลักในการแก้โจทย์ของไหล}

\begin{tcolorbox}[title=ขั้นตอนแก้โจทย์ของไหล]
\begin{enumerate}
    \item ดูก่อนว่าเป็นของไหลอยู่นิ่งหรือกำลังไหล
    \item ถ้าอยู่นิ่ง ใช้
    \[
    P=P_0+\rho gh
    \]
    และหลักความดันเท่ากันที่ระดับเดียวกัน
    \item ถ้ามีวัตถุลอยหรือจม ใช้
    \[
    B=\rho_f gV_{\text{disp}}
    \]
    \item ถ้ากำลังไหลและของไหลอัดตัวไม่ได้ ใช้
    \[
    A_1v_1=A_2v_2
    \]
    \item ถ้าเกี่ยวกับความดัน ความสูง และอัตราเร็วในของไหลกำลังไหล ใช้
    \[
    P+\frac{1}{2}\rho v^2+\rho gy=\text{constant}
    \]
    \item ถ้ามีแรงต้านหรือความตึงผิว ใช้สมการที่โจทย์กำหนดหรือสูตรเฉพาะของหัวข้อนั้น
\end{enumerate}
\end{tcolorbox}

% =========================================================
\section{ชุดโจทย์รวมท้ายบท: ระดับผสมสำหรับ สอวน. ค่าย 1}

\subsection*{ระดับพื้นฐาน}

\begin{enumerate}
    \item จงหาความดันเกจที่ความลึก $3\,\si{m}$ ในน้ำ ใช้ $g=10\,\si{m/s^2}$
    \item วัตถุปริมาตร $5\times 10^{-4}\,\si{m^3}$ จมอยู่ในน้ำทั้งหมด จงหาแรงพยุง
    \item น้ำไหลผ่านท่อพื้นที่หน้าตัด $4\,\si{cm^2}$ ด้วยอัตราเร็ว $3\,\si{m/s}$ จงหาอัตราการไหล
    \item ถังเปิดมีรูเล็กที่ลึกจากผิวน้ำ $0.80\,\si{m}$ ใช้ $g=10\,\si{m/s^2}$ จงหาอัตราเร็วของน้ำที่พุ่งออก
\end{enumerate}

\subsection*{ระดับกลาง}

\begin{enumerate}
    \setcounter{enumi}{4}
    \item ก้อนไม้ความหนาแน่น $600\,\si{kg/m^3}$ ลอยในน้ำ จงหาสัดส่วนปริมาตรที่จมน้ำ
    \item หลอดรูปตัวยูมีของเหลวสองชนิด ความหนาแน่น $\rho_1$ และ $\rho_2$ อยู่สมดุลกัน จงเขียนหลักการตั้งสมการความดันที่ระดับเดียวกัน
    \item ลูกปิงปองมวล $m$ ปริมาตร $V$ อยู่ใต้น้ำ ถ้าไม่คิดแรงต้านน้ำ จงหาแรงลัพธ์ทันทีหลังปล่อย
    \item ถังเปิดมีพื้นที่หน้าตัด $A$ และรูพื้นที่ $a$ ที่ก้นถัง ความสูงน้ำคือ $h$ จงพิสูจน์ว่า
    \[
    v=
    \sqrt{
    \frac{2gh}{1-\left(\frac{a}{A}\right)^2}
    }
    \]
\end{enumerate}

\subsection*{ระดับยาก}

\begin{enumerate}
    \setcounter{enumi}{8}
    \item วัตถุลอยในของเหลวชนิดหนึ่งโดยจม $70\%$ ของปริมาตร ถ้าความหนาแน่นของของเหลวคือ $\rho_f$ จงหาความหนาแน่นของวัตถุ
    \item น้ำไหลในท่อแนวนอนจากพื้นที่หน้าตัด $A_1$ ไป $A_2$ โดย $A_2<A_1$ จงเปรียบเทียบความดันที่สองตำแหน่ง
    \item ฟองสบู่รัศมี $R$ มีความตึงผิว $\gamma$ จงหาความดันในฟองที่มากกว่าความดันภายนอก
    \item วัตถุในของไหลมีแรงต้าน $F_D=Cv$ และแรงขับคงตัว $F_0$ จงหาความเร็วปลาย
\end{enumerate}

% =========================================================
\section{เฉลยย่อท้ายบท}

\begin{enumerate}
    \item
    \[
    P_{\text{gauge}}=\rho gh=(1000)(10)(3)=3.0\times10^4\,\si{Pa}
    \]

    \item
    \[
    B=\rho gV=(1000)(10)(5\times10^{-4})=5\,\si{N}
    \]

    \item
    \[
    Q=Av
    \]
    \[
    A=4\,\si{cm^2}=4\times10^{-4}\,\si{m^2}
    \]
    \[
    Q=(4\times10^{-4})(3)=1.2\times10^{-3}\,\si{m^3/s}
    \]

    \item
    \[
    v=\sqrt{2gh}=\sqrt{2(10)(0.80)}=4\,\si{m/s}
    \]

    \item
    \[
    \frac{V_{\text{disp}}}{V_{\text{object}}}
    =
    \frac{\rho_{\text{wood}}}{\rho_{\text{water}}}
    =
    \frac{600}{1000}=0.60
    \]

    \item
    เลือกระดับเดียวกันในของเหลวที่ต่อถึงกัน แล้วเขียน
    \[
    P_{\text{left}}=P_{\text{right}}
    \]
    โดยเพิ่มพจน์ $\rho gh$ ทุกครั้งที่ลงลึกในของเหลว

    \item
    แรงพยุงขึ้นคือ
    \[
    B=\rho_{\text{water}}gV
    \]
    น้ำหนักลงคือ
    \[
    mg
    \]
    ดังนั้นแรงลัพธ์ขึ้นคือ
    \[
    F_{\text{net}}=\rho_{\text{water}}gV-mg
    \]

    \item
    จากความต่อเนื่อง
    \[
    Av_s=av
    \]
    \[
    v_s=\frac{a}{A}v
    \]
    จากแบร์นูลลี
    \[
    gh+\frac{1}{2}v_s^2=\frac{1}{2}v^2
    \]
    แทน $v_s$
    \[
    gh+\frac{1}{2}\left(\frac{a}{A}\right)^2v^2=\frac{1}{2}v^2
    \]
    จึงได้
    \[
    v=
    \sqrt{
    \frac{2gh}{1-\left(\frac{a}{A}\right)^2}
    }
    \]

    \item
    \[
    \frac{V_{\text{disp}}}{V_{\text{object}}}
    =
    \frac{\rho_{\text{object}}}{\rho_f}
    \]
    ดังนั้น
    \[
    \rho_{\text{object}}=0.70\rho_f
    \]

    \item
    จากความต่อเนื่อง
    \[
    A_1v_1=A_2v_2
    \]
    เมื่อ $A_2<A_1$ จะได้
    \[
    v_2>v_1
    \]
    ในท่อแนวนอน จากแบร์นูลลี
    \[
    P+\frac{1}{2}\rho v^2=\text{constant}
    \]
    ดังนั้นตำแหน่งที่ไหลเร็วกว่า มีความดันต่ำกว่า

    \item
    ฟองสบู่มีสองผิว
    \[
    P_i-P_o=\frac{4\gamma}{R}
    \]

    \item
    ที่ความเร็วปลาย
    \[
    \sum F=0
    \]
    \[
    F_0-Cv_t=0
    \]
    ดังนั้น
    \[
    v_t=\frac{F_0}{C}
    \]
\end{enumerate}

% =========================================================
\section{สรุปสูตรสำคัญ}

\begin{tcolorbox}[title=สรุปบท: กลศาสตร์ของไหล]
ความหนาแน่น

\[
\rho=\frac{m}{V}
\]

ความดัน

\[
P=\frac{F}{A}
\]

ความดันในของเหลวอยู่นิ่ง

\[
P=P_0+\rho gh
\]

ความดันเกจ

\[
P_{\text{gauge}}=\rho gh
\]

ความดันสัมบูรณ์

\[
P_{\text{abs}}=P_{\text{atm}}+P_{\text{gauge}}
\]

หลักของปาสคาล

\[
\frac{F_1}{A_1}=\frac{F_2}{A_2}
\]

แรงพยุง

\[
B=\rho_f gV_{\text{disp}}
\]

วัตถุลอยนิ่ง

\[
\frac{V_{\text{disp}}}{V_{\text{object}}}
=
\frac{\rho_{\text{object}}}{\rho_f}
\]
\end{tcolorbox}

\begin{tcolorbox}[title=สรุปบท: กลศาสตร์ของไหล ต่อ]
อัตราการไหล

\[
Q=Av
\]

สมการความต่อเนื่อง

\[
A_1v_1=A_2v_2
\]

สมการแบร์นูลลี

\[
P+\frac{1}{2}\rho v^2+\rho gy=\text{constant}
\]

กฎของตอร์ริเชลลี

\[
v=\sqrt{2gh}
\]

รูที่ก้นถังเมื่อพื้นที่ถังไม่ใหญ่มาก

\[
v=
\sqrt{
\frac{2gh}
{1-\left(\frac{a}{A}\right)^2}
}
\]

แรงต้านแบบเส้นตรง

\[
F_D=Cv
\]

แรงต้านแบบกำลังสอง

\[
F_D=kv^2A
\]

ฟองสบู่

\[
P_i-P_o=\frac{4\gamma}{R}
\]
\end{tcolorbox}

% =========================================================
\section{ข้อสอบจริง สอวน. ที่ใช้ความรู้เรื่องของไหล}



% =========================================================
\subsection{ปี 2560 ข้อ 1: แรงต้านในของไหลและมิติของค่าคงตัว}

\vspace{1.5em}

\IfFileExists{img/posn60q1.png}{
\begin{center}
    \includegraphics[width=1\linewidth]{img/posn60q1.png}
\end{center}
}{
\begin{tcolorbox}[title=ตำแหน่งภาพที่แนะนำ]
ให้ตัดภาพข้อสอบปี 2560 ข้อ 1 แล้วบันทึกเป็น

\[
\texttt{img/posn60q1.png}
\]
\end{tcolorbox}
}

\begin{examplebox}{เฉลยละเอียด}
โจทย์ให้แรงต้านในของไหลเป็น

\[
F=kv^2A
\]

ต้องการหาว่า $k$ มีมิติเป็นปริมาณใด

เริ่มจากจัดรูป

\[
k=\frac{F}{v^2A}
\]

พิจารณาหน่วย

\[
[F]=\si{N}=\si{kg.m/s^2}
\]

\[
[v^2]=\si{m^2/s^2}
\]

\[
[A]=\si{m^2}
\]

ดังนั้น

\[
[k]
=
\frac{\si{kg.m/s^2}}
{\left(\si{m^2/s^2}\right)\left(\si{m^2}\right)}
\]

\[
[k]=\si{kg/m^3}
\]

ซึ่งเป็นหน่วยของความหนาแน่น

\[
\boxed{k \text{ เป็นความหนาแน่น}}
\]
\end{examplebox}

% =========================================================
\subsection{ปี 2560 ข้อ 4: ลูกปิงปองลอยขึ้นในน้ำและความเร็วปลาย}

\vspace{1.5em}

\IfFileExists{img/posn60q4.png}{
\begin{center}
    \includegraphics[width=1\linewidth]{img/posn60q4.png}
\end{center}
}{
\begin{tcolorbox}[title=ตำแหน่งภาพที่แนะนำ]
ให้ตัดภาพข้อสอบปี 2560 ข้อ 4 แล้วบันทึกเป็น

\[
\texttt{img/posn60q4.png}
\]
\end{tcolorbox}
}

\begin{examplebox}{เฉลยละเอียด}
ลูกปิงปองมวล $m$ ปริมาตร $V$ อยู่ใต้น้ำความหนาแน่น $\rho$ และมีแรงต้านน้ำขณะลอยขึ้น

\[
F_D=Cv
\]

แรงที่กระทำมีสามแรง

แรงพยุงขึ้น

\[
B=\rho gV
\]

น้ำหนักลง

\[
mg
\]

แรงต้านลง เพราะลูกปิงปองกำลังเคลื่อนที่ขึ้น

\[
F_D=Cv
\]

เมื่อถึงความเร็วคงตัว

\[
\sum F=0
\]

เลือกทิศขึ้นเป็นบวก

\[
\rho gV-mg-Cv=0
\]

ดังนั้น

\[
Cv=\rho gV-mg
\]

\[
v=\frac{\rho gV-mg}{C}
\]

สรุป

\[
\boxed{v=\frac{\rho gV-mg}{C}}
\]
\end{examplebox}

% =========================================================
\subsection{ปี 2560 ข้อ 16: ทรงกลมลอยน้ำและอัตราส่วนความหนาแน่น}

\vspace{1.5em}

\IfFileExists{img/posn60q16.png}{
\begin{center}
    \includegraphics[width=1\linewidth]{img/posn60q16.png}
\end{center}
}{
\begin{tcolorbox}[title=ตำแหน่งภาพที่แนะนำ]
ให้ตัดภาพข้อสอบปี 2560 ข้อ 16 แล้วบันทึกเป็น

\[
\texttt{img/posn60q16.png}
\]
\end{tcolorbox}
}

\begin{examplebox}{เฉลยละเอียด}
สำหรับวัตถุลอยนิ่ง

\[
\frac{\rho_{\text{object}}}{\rho_{\text{water}}}
=
\frac{V_{\text{sub}}}{V_{\text{total}}}
\]

\textbf{ทรงกลมแรก}

ปริมาตรรวมคือ

\[
V_1=\frac{4}{3}\pi a^3
\]

โจทย์ให้ปริมาตรส่วนที่พ้นน้ำเป็น

\[
V_{\text{above}}=\frac{\pi a^3}{3}
\]

ดังนั้นปริมาตรส่วนที่จมน้ำคือ

\[
V_{\text{sub},1}
=
\frac{4}{3}\pi a^3-\frac{1}{3}\pi a^3
\]

\[
V_{\text{sub},1}=\pi a^3
\]

จึงได้

\[
\frac{\rho_1}{\rho_w}
=
\frac{\pi a^3}{\frac{4}{3}\pi a^3}
=
\frac{3}{4}
\]

\textbf{ทรงกลมที่สอง}

ปริมาตรรวมคือ

\[
V_2=\frac{4}{3}\pi b^3
\]

โจทย์ให้ปริมาตรส่วนที่จมน้ำเป็น

\[
V_{\text{sub},2}=\frac{8\pi b^3}{9}
\]

ดังนั้น

\[
\frac{\rho_2}{\rho_w}
=
\frac{\frac{8\pi b^3}{9}}{\frac{4}{3}\pi b^3}
\]

\[
\frac{\rho_2}{\rho_w}
=
\frac{2}{3}
\]

อัตราส่วนความหนาแน่นคือ

\[
\rho_1:\rho_2
=
\frac{3}{4}:\frac{2}{3}
\]

คูณด้วย $12$

\[
\rho_1:\rho_2=9:8
\]

ดังนั้น

\[
\boxed{9:8}
\]
\end{examplebox}

% =========================================================
\subsection{ปี 2560 ข้อ 17: หลอดปรอทและความดันเหนือผิวปรอท}

\vspace{1.5em}

\IfFileExists{img/posn60q17.png}{
\begin{center}
    \includegraphics[width=1\linewidth]{img/posn60q17.png}
\end{center}
}{
\begin{tcolorbox}[title=ตำแหน่งภาพที่แนะนำ]
ให้ตัดภาพข้อสอบปี 2560 ข้อ 17 แล้วบันทึกเป็น

\[
\texttt{img/posn60q17.png}
\]
\end{tcolorbox}
}

\begin{examplebox}{เฉลยละเอียด}
ความดันบรรยากาศคือ

\[
750\,\si{mmHg}
\]

ผิวปรอทในหลอดสูงกว่าผิวปรอทในอ่าง

\[
10\,\si{mm}
\]

ให้ความดันเหนือผิวปรอทในหลอดเป็น

\[
P
\]

ที่ระดับเดียวกันในปรอท ความดันต้องเท่ากัน ดังนั้น

\[
P+10=750
\]

จึงได้

\[
P=740\,\si{mmHg}
\]

ดังนั้น

\[
\boxed{740\,\si{mmHg}}
\]
\end{examplebox}

% =========================================================
\subsection{ปี 2560 ตอนที่ 2 ข้อ 8: ถังน้ำมีรูที่ก้นและอัตราการไหล}

\vspace{1.5em}

\IfFileExists{img/posn60s2q8.png}{
\begin{center}
    \includegraphics[width=1\linewidth]{img/posn60s2q8.png}
\end{center}
}{
\begin{tcolorbox}[title=ตำแหน่งภาพที่แนะนำ]
ให้ตัดภาพข้อสอบปี 2560 ตอนที่ 2 ข้อ 8 แล้วบันทึกเป็น

\[
\texttt{img/posn60s2q8.png}
\]
\end{tcolorbox}
}

\begin{examplebox}{เฉลยแนวคิด}
โจทย์ให้ถังทรงกระบอกเปิดฝา รัศมีถัง

\[
R=1.0\,\si{m}
\]

ดังนั้นพื้นที่หน้าตัดถังคือ

\[
A=\pi R^2=\pi\,\si{m^2}
\]

อัตราการไหลออกจากรูคือ

\[
Q
\]

อัตราเร็วของผิวน้ำในถังคือ

\[
v_{\text{surface}}=\frac{Q}{A}
\]

ถ้าโจทย์ให้

\[
Q=8\pi\,\si{m^3/s}
\]

จะได้

\[
v_{\text{surface}}=\frac{8\pi}{\pi}=8\,\si{m/s}
\]

แปลงเป็นเซนติเมตรต่อวินาที

\[
8\,\si{m/s}=800\,\si{cm/s}
\]

ดังนั้น

\[
\boxed{v_{\text{surface}}=800\,\si{cm/s}}
\]
\end{examplebox}

% =========================================================
\subsection{ปี 2561 ข้อ 16: ของเหลวสามชนิดในหลอดรูปตัวยู}

\vspace{1.5em}

\IfFileExists{img/posn61q16.png}{
\begin{center}
    \includegraphics[width=1\linewidth]{img/posn61q16.png}
\end{center}
}{
\begin{tcolorbox}[title=ตำแหน่งภาพที่แนะนำ]
ให้ตัดภาพข้อสอบปี 2561 ข้อ 16 แล้วบันทึกเป็น

\[
\texttt{img/posn61q16.png}
\]
\end{tcolorbox}
}

\begin{examplebox}{เฉลยละเอียด}
โจทย์เป็นหลอดรูปตัวยูที่มีของเหลวหลายชั้น หลักสำคัญคือเลือกจุดสองจุดในของเหลวชนิดเดียวกันที่อยู่ระดับเดียวกัน แล้วใช้

\[
P_{\text{left}}=P_{\text{right}}
\]

จากรูป ผิวของเหลวทั้งสองด้านอยู่ระดับเดียวกัน และให้ระยะด้านซ้ายเป็น

\[
B=12\,\si{cm}
\]

ต้องการหาระยะ

\[
A
\]

ให้ของเหลวด้านซ้ายมีความหนาแน่น $900\,\si{kg/m^3}$ ด้านขวา $800\,\si{kg/m^3}$ และชั้นล่างสุด $1100\,\si{kg/m^3}$

ตั้งสมการความดันระหว่างรอยต่อทั้งสองโดยคำนึงถึงชั้นของเหลวล่างที่คั่นอยู่

\[
900B=800A+1100(B-A)
\]

แทน

\[
B=12
\]

\[
900(12)=800A+1100(12-A)
\]

\[
10800=800A+13200-1100A
\]

\[
300A=2400
\]

ดังนั้น

\[
A=8\,\si{cm}
\]

สรุป

\[
\boxed{8\,\si{cm}}
\]
\end{examplebox}

% =========================================================
\subsection{ปี 2561 ตอนที่ 2 ข้อ 8: แท่งไม้บางจมน้ำบางส่วนและสมดุลทอร์ก}

\vspace{1.5em}

\IfFileExists{img/posn61s2q8.png}{
\begin{center}
    \includegraphics[width=1\linewidth]{img/posn61s2q8.png}
\end{center}
}{
\begin{tcolorbox}[title=ตำแหน่งภาพที่แนะนำ]
ให้ตัดภาพข้อสอบปี 2561 ตอนที่ 2 ข้อ 8 แล้วบันทึกเป็น

\[
\texttt{img/posn61s2q8.png}
\]
\end{tcolorbox}
}

\begin{examplebox}{เฉลยละเอียด}
แท่งไม้บางยาว

\[
L=60\,\si{cm}
\]

ปลายหนึ่งยึดกับแกนหมุน และปลายอีกด้านจมน้ำ โดยส่วนที่จมน้ำยาว

\[
\ell=20\,\si{cm}
\]

ให้พื้นที่หน้าตัดของแท่งเป็น $A$ ความหนาแน่นของแท่งเป็น $\rho_{\text{wood}}$ และความหนาแน่นของน้ำเป็น $\rho_w$

น้ำหนักแท่งคือ

\[
W=\rho_{\text{wood}}gAL
\]

แรงนี้กระทำที่กึ่งกลางแท่ง ห่างจากจุดหมุน

\[
\frac{L}{2}=30\,\si{cm}
\]

แรงพยุงเกิดกับส่วนที่จมน้ำ

\[
B=\rho_w gA\ell
\]

และกระทำที่กึ่งกลางของส่วนที่จมน้ำ ซึ่งอยู่ห่างจากจุดหมุน

\[
L-\frac{\ell}{2}
=
60-10=50\,\si{cm}
\]

สมดุลทอร์กรอบจุดหมุน

\[
B(50)=W(30)
\]

แทนค่า

\[
(\rho_w gA\ell)(50)
=
(\rho_{\text{wood}}gAL)(30)
\]

ตัด $gA$ ออก

\[
\rho_w(20)(50)=\rho_{\text{wood}}(60)(30)
\]

ดังนั้น

\[
\frac{\rho_w}{\rho_{\text{wood}}}
=
\frac{60(30)}{20(50)}
=
\frac{1800}{1000}
\]

\[
\frac{\rho_w}{\rho_{\text{wood}}}=\frac{9}{5}
\]

ดังนั้นน้ำมีความหนาแน่นเป็น

\[
\boxed{\frac{9}{5}}
\]

เท่าของแท่งไม้
\end{examplebox}

% =========================================================
\subsection{ปี 2562 ข้อ 3: กาลักน้ำและความดันที่ยอดท่อ}

\vspace{1.5em}

\IfFileExists{img/posn62q3.png}{
\begin{center}
    \includegraphics[width=1\linewidth]{img/posn62q3.png}
\end{center}
}{
\begin{tcolorbox}[title=ตำแหน่งภาพที่แนะนำ]
ให้ตัดภาพข้อสอบปี 2562 ข้อ 3 แล้วบันทึกเป็น

\[
\texttt{img/posn62q3.png}
\]
\end{tcolorbox}
}

\begin{examplebox}{เฉลยละเอียด}
ที่ยอดท่อกาลักน้ำ ความดันต้องไม่ต่ำกว่าสุญญากาศ

ให้ความดันบรรยากาศเป็น

\[
P_a
\]

ความหนาแน่นของน้ำเป็น

\[
\rho
\]

และยอดท่อสูงกว่าผิวน้ำเป็น

\[
h
\]

จากความดันในของเหลว

\[
P_{\text{top}}=P_a-\rho gh
\]

เพื่อให้กาลักน้ำทำงานได้ในแบบจำลองพื้นฐาน ต้องมี

\[
P_{\text{top}}>0
\]

ดังนั้น

\[
P_a-\rho gh>0
\]

\[
h<\frac{P_a}{\rho g}
\]

สรุป

\[
\boxed{h<\frac{P_a}{\rho g}}
\]
\end{examplebox}

% =========================================================
\subsection{ปี 2562 ข้อ 4: อัตราเร็วของน้ำที่พุ่งจากก้นถัง}

\vspace{1.5em}

\IfFileExists{img/posn62q4.png}{
\begin{center}
    \includegraphics[width=1\linewidth]{img/posn62q4.png}
\end{center}
}{
\begin{tcolorbox}[title=ตำแหน่งภาพที่แนะนำ]
ให้ตัดภาพข้อสอบปี 2562 ข้อ 4 แล้วบันทึกเป็น

\[
\texttt{img/posn62q4.png}
\]
\end{tcolorbox}
}

\begin{examplebox}{เฉลยละเอียด}
ให้พื้นที่หน้าตัดถังเป็น

\[
A
\]

พื้นที่รูเป็น

\[
a
\]

และความสูงของน้ำเหนือรูเป็น

\[
h
\]

จากความต่อเนื่อง

\[
Av_s=av
\]

ดังนั้น

\[
v_s=\frac{a}{A}v
\]

ใช้แบร์นูลลีระหว่างผิวน้ำกับรู

\[
\rho gh+\frac{1}{2}\rho v_s^2
=
\frac{1}{2}\rho v^2
\]

ตัด $\rho$ ออก

\[
gh+\frac{1}{2}v_s^2=\frac{1}{2}v^2
\]

แทน

\[
v_s=\frac{a}{A}v
\]

\[
gh+\frac{1}{2}\left(\frac{a}{A}\right)^2v^2
=
\frac{1}{2}v^2
\]

จัดรูป

\[
2gh
=
v^2
\left[
1-\left(\frac{a}{A}\right)^2
\right]
\]

ดังนั้น

\[
\boxed{
v=
\sqrt{
\frac{2gh}
{1-\left(\frac{a}{A}\right)^2}
}
}
\]
\end{examplebox}

% =========================================================
\subsection{ปี 2564 ข้อ 6: ลูกปิงปองถูกปล่อยจากก้นถ้วยน้ำ}

\vspace{1.5em}

\IfFileExists{img/posn64q6.png}{
\begin{center}
    \includegraphics[width=1\linewidth]{img/posn64q6.png}
\end{center}
}{
\begin{tcolorbox}[title=ตำแหน่งภาพที่แนะนำ]
ให้ตัดภาพข้อสอบปี 2564 ข้อ 6 แล้วบันทึกเป็น

\[
\texttt{img/posn64q6.png}
\]
\end{tcolorbox}
}

\begin{examplebox}{เฉลยละเอียด}
ลูกปิงปองมวล $m$ ปริมาตร $V$ ถูกปล่อยจากก้นถ้วยน้ำลึก $h$ ไม่คิดแรงต้านน้ำและอากาศ

ระหว่างที่อยู่ใต้น้ำ แรงขึ้นคือแรงพยุง

\[
B=\rho gV
\]

แรงลงคือน้ำหนัก

\[
mg
\]

แรงลัพธ์ขึ้นคือ

\[
B-mg=\rho gV-mg
\]

งานสุทธิที่ทำระหว่างลอยขึ้นในน้ำเป็นระยะ $h$ คือ

\[
W_{\text{net}}=(\rho gV-mg)h
\]

งานนี้กลายเป็นพลังงานจลน์ที่ผิวน้ำ

\[
K=(\rho gV-mg)h
\]

หลังพ้นน้ำ แรงพยุงหายไป เหลือแต่น้ำหนัก ทำให้ลูกปิงปองขึ้นต่อได้สูง $H$ โดย

\[
mgH=K
\]

ดังนั้น

\[
mgH=(\rho gV-mg)h
\]

ตัด $g$ ออก

\[
mH=(\rho V-m)h
\]

\[
H=\left(\frac{\rho V}{m}-1\right)h
\]

สรุป

\[
\boxed{H=\left(\frac{\rho V}{m}-1\right)h}
\]
\end{examplebox}

% =========================================================
\subsection{ปี 2566 ข้อ 7: จำนวนโมเลกุลทั้งหมดในบรรยากาศจากความดัน}

\vspace{1.5em}

\IfFileExists{img/posn66q7.png}{
\begin{center}
    \includegraphics[width=1\linewidth]{img/posn66q7.png}
\end{center}
}{
\begin{tcolorbox}[title=ตำแหน่งภาพที่แนะนำ]
ให้ตัดภาพข้อสอบปี 2566 ข้อ 7 แล้วบันทึกเป็น

\[
\texttt{img/posn66q7.png}
\]
\end{tcolorbox}
}

\begin{examplebox}{เฉลยละเอียด}
แนวคิดคือ ความดันบรรยากาศที่ผิวโลกเกิดจากน้ำหนักรวมของบรรยากาศต่อพื้นที่ผิวโลก

พื้นที่ผิวโลกประมาณ

\[
A=4\pi R^2
\]

แรงดันรวมจากบรรยากาศคือ

\[
F=PA
\]

แรงนี้เท่ากับน้ำหนักของบรรยากาศทั้งหมด

\[
F=Nmg
\]

ดังนั้น

\[
P(4\pi R^2)=Nmg
\]

จัดรูป

\[
N=\frac{4\pi R^2P}{mg}
\]

แทนค่าประมาณ

\[
R\approx 6\times10^6\,\si{m}
\]

\[
P\approx 10^5\,\si{Pa}
\]

\[
m\approx 5\times10^{-26}\,\si{kg}
\]

\[
g\approx 10\,\si{m/s^2}
\]

จะได้

\[
N\sim
\frac{4\pi(6\times10^6)^2(10^5)}
{(5\times10^{-26})(10)}
\]

ประมาณลำดับขนาด

\[
N\sim 10^{44}
\]

ดังนั้น

\[
\boxed{N\sim 10^{44}}
\]
\end{examplebox}

% =========================================================
\subsection{ปี 2566 ข้อ 11: น้ำแข็งลอยและระดับน้ำหลังละลาย}

\vspace{1.5em}

\IfFileExists{img/posn66q11.png}{
\begin{center}
    \includegraphics[width=1\linewidth]{img/posn66q11.png}
\end{center}
}{
\begin{tcolorbox}[title=ตำแหน่งภาพที่แนะนำ]
ให้ตัดภาพข้อสอบปี 2566 ข้อ 11 แล้วบันทึกเป็น

\[
\texttt{img/posn66q11.png}
\]
\end{tcolorbox}
}

\begin{examplebox}{เฉลยละเอียด}
น้ำแข็งลอยนิ่งในน้ำ ดังนั้นแรงพยุงเท่ากับน้ำหนักน้ำแข็ง

\[
\rho_w gV_{\text{sub}}=\rho_{\text{ice}}gV_{\text{ice}}
\]

ดังนั้นปริมาตรน้ำที่ถูกแทนที่ตอนน้ำแข็งยังไม่ละลายคือ

\[
V_{\text{disp}}=
\frac{\rho_{\text{ice}}}{\rho_w}V_{\text{ice}}
\]

เมื่อน้ำแข็งละลายหมด มวลเท่าเดิม กลายเป็นน้ำปริมาตร

\[
V_{\text{melt}}=
\frac{m_{\text{ice}}}{\rho_w}
=
\frac{\rho_{\text{ice}}V_{\text{ice}}}{\rho_w}
\]

ซึ่งเท่ากับ

\[
V_{\text{disp}}
\]

ดังนั้นปริมาตรน้ำที่เพิ่มหลังละลายเท่ากับปริมาตรน้ำที่เคยถูกแทนที่พอดี

ระดับน้ำจึง

\[
\boxed{\text{ไม่เปลี่ยนแปลง}}
\]
\end{examplebox}

% =========================================================
\subsection{ปี 2566 ข้อ 14: ผิวน้ำในถ้วยที่หมุน}

\vspace{1.5em}

\IfFileExists{img/posn66q14.png}{
\begin{center}
    \includegraphics[width=1\linewidth]{img/posn66q14.png}
\end{center}
}{
\begin{tcolorbox}[title=ตำแหน่งภาพที่แนะนำ]
ให้ตัดภาพข้อสอบปี 2566 ข้อ 14 แล้วบันทึกเป็น

\[
\texttt{img/posn66q14.png}
\]
\end{tcolorbox}
}

\begin{examplebox}{เฉลยละเอียด}
ของเหลวหมุนพร้อมภาชนะด้วยอัตราเร็วเชิงมุม

\[
\omega
\]

ที่ระยะ $x$ จากแกนหมุน ของเหลวต้องมีความเร่งสู่ศูนย์กลาง

\[
a_c=\omega^2x
\]

ในกรอบที่หมุนพร้อมของเหลว ผิวน้ำต้องตั้งฉากกับความเร่งลัพธ์เสมือน ซึ่งมาจาก

\[
g
\]

ลงด้านล่าง และ

\[
\omega^2x
\]

ออกด้านข้าง

ดังนั้นความชันของผิวน้ำเป็น

\[
\frac{dy}{dx}=\frac{\omega^2x}{g}
\]

อินทิเกรต

\[
dy=\frac{\omega^2}{g}x\,dx
\]

\[
y=\frac{\omega^2}{2g}x^2+C
\]

ถ้าให้ที่แกนกลางมีระดับ

\[
y(0)=h
\]

จะได้

\[
C=h
\]

ดังนั้น

\[
\boxed{
y=\frac{\omega^2}{2g}x^2+h
}
\]
\end{examplebox}

% =========================================================
\subsection{ปี 2567 ข้อ 13: ฟองสบู่และความตึงผิว}

\vspace{1.5em}

\IfFileExists{img/posn67q13.png}{
\begin{center}
    \includegraphics[width=1\linewidth]{img/posn67q13.png}
\end{center}
}{
\begin{tcolorbox}[title=ตำแหน่งภาพที่แนะนำ]
ให้ตัดภาพข้อสอบปี 2567 ข้อ 13 แล้วบันทึกเป็น

\[
\texttt{img/posn67q13.png}
\]
\end{tcolorbox}
}

\begin{examplebox}{เฉลยละเอียด}
ฟองสบู่มีผิวสองด้าน คือผิวด้านนอกและผิวด้านใน

สำหรับผิวทรงกลมหนึ่งผิว ความดันด้านในมากกว่าด้านนอกเป็น

\[
\Delta P=\frac{2\gamma}{R}
\]

แต่ฟองสบู่มีสองผิว จึงได้ความดันเกินเป็นสองเท่า

\[
\Delta P=\frac{4\gamma}{R}
\]

ดังนั้น

\[
P_i-P_o=\frac{4\gamma}{R}
\]

หรือ

\[
\boxed{
P_i=P_o+\frac{4\gamma}{R}
}
\]
\end{examplebox}